\subsection{Rank, norm, and even the entire interaction spectrum are
insufficient}\label{proof:local_interactions:rank-norm-and-even-the-entire-interaction-spectrum-are-insufficient}

The following comparison uses zero-diagonal matrices throughout, so it
is not an artifact of adding diagonal constants to the Hamiltonian.

Let n be even. Define \(J_{\mathrm{loc}}\) to have \(J_{12}=J_{21}=1\)
and all other entries zero. Define \(J_{\mathrm{deloc}}\) on two equal
parts \(L,R\) of size \(N=n/2\) by

\[
(J_{\rm deloc})_{ij}=\begin{cases}1/N,&i\in L,j\in R\text{ or vice versa},\\0,&\text{otherwise}.\end{cases}
\]

Both matrices have exactly the same full spectrum

\[
\{1,-1,\underbrace{0,\ldots,0}_{n-2}\},
\]

and therefore the same rank two, operator norm one, and Frobenius norm
\(\sqrt{2}\). Both have maximum absolute row sum one. For
\(J_{\mathrm{loc}}\) the two nonzero eigenvectors are supported on two
sites. For \(J_{\mathrm{deloc}}\) they are constant or staggered over
all n sites.

The \(J_{\mathrm{loc}}\) Gibbs law is one correlated pair and \(n-2\)
independent fair spins, so the bounded-block theorem gives AA at every
fixed finite \(\beta\).

For \(J_{\mathrm{deloc}}\), take S with a sites in L and b in R, and
write \(k=a+b\le n/2\). Its weighted boundary is

\[
B_{J_{\rm deloc}}(S)=k-\frac{4ab}{n}
\ge k-\frac{k^2}{n}\ge k/2.
\]

Thus \(\kappa =1/2\) in Theorem L2. Taking \(\alpha =1/4\) proves
failure for every fixed \(\beta >4\log 2\), in particular at
\(\beta =6\). The law has a distributed magnetization witness, while the
isospectral localized matrix does not.

\textbf{Conclusion.} Neither fixed interaction rank, bounded operator
norm, row normalization, nor the entire list of eigenvalues determines
the AA outcome of an individual family. Structural criteria using
spectral information must also control where the relevant eigenvectors
live or otherwise detect collective order. This does not refute a
small-enough spectral-norm sufficient condition; the counterexample
operates at a specified larger value of \(\beta \lVert J\rVert\).

The isospectral comparison can also have an extensive energy range on
both sides. For n divisible by four, use the preceding comparison on the
first \(n/2\) vertices and append the same disjoint unit-weight matching
on the other \(n/2\) vertices. The two full matrices remain exactly
isospectral and row-normalized, and both Hamiltonians now have energy
range \(\Theta (n)\). The positive member still has components of size
at most two; the negative member retains the ordered component of size
\(n/2\), whose normalized magnetization is a degree-one AA
counterexample. Fixed rank and extensivity are different constraints;
this variant preserves the isospectral obstruction while giving up fixed
total rank.

Diagonal convention matters in any proposed spectral criterion:
replacing J by \(J+cI\) changes H only by the constant \(-cn/2\) and
leaves the Gibbs law unchanged. Therefore positivity of J alone, without
a fixed diagonal convention or an optimization over diagonal shifts,
imposes no substantive restriction on the classical Gibbs distribution.
